\section{Schubert coefficients of sparse paving matroids}
In this section we compute the Schubert coefficients of sparse paving matroids. The observation that allows us to do so is the following.
\begin{proposition}
    \label{prop:VolDeg}
    For $x \in G(r,n)$ let $s = \dim(\T) = n - \kappa(M_x)$. The following equality holds
    \begin{equation}
        \label{eq:impLem}
        \deg([\T]\sigma_{(1)}^s) = \Vol_s(P(M_x)).
    \end{equation}
    This gives the linear relation among the Schubert coefficients of $M_x$ 
    \begin{equation}
        \label{eq:linRel}
        \sum_\lambda c_{\lambda^c} d_\lambda(M_x) = \Vol_s(P(M_x)),
    \end{equation}
    where $c_{\lambda^c}$ are the integer coefficients found in the expansion of $\sigma_{(1)}^s$. That is,  $c_{\lambda^c} = \deg(\sigma_{(1)}^s \sigma_{\lambda}).$ Here the sum is taken over all partitions $\lambda \subseteq r\times(n-r)$ with $|\lambda| = r(n-r) - s$.
\end{proposition}
\begin{proof}
    There are two ways of computing the degree $D$ of $\T$ in $\PP^{\binom{n}{r}-1}$. First, since $\T$ is a subvariety of $G(r,n)$, $D = \deg([\T]\sigma_{(1)}^s)$. Second, since $\T$ is isomorphic to the projective toric variety obtained from $P(M_x)$, $D = Vol_s(P(M_x))$. Expanding $[\T]$ and $\sigma_{(1)}^s$ gives
    \begin{align*}
        [\T] &= \sum_\lambda d_\lambda(M_x)\sigma_\lambda \text{ and }\\
        \sigma_{(1)}^s &= \sum_\lambda c_{\lambda^c} \sigma_{\lambda^c}.
    \end{align*}
    Now by the complimentary dimension formula the left hand side of \cref{eq:impLem} becomes the left hand side of \cref{eq:linRel}.
\end{proof}

The proof above only works for matroids representable over $\CC$, but we can extend the result to arbitrary matroids using the fact that valuative matroid invariants are determined by their values on Schubert matroids.
\begin{corollary}
    Let $M$ be a rank $r$ matroid on $[n]$ and $s = n - \kappa(M)$ then 
    \begin{equation*}
        \deg(\Sc(M)\sigma_{(1)}^s) = Vol_s(P(M)).
    \end{equation*}
    This gives the linear relation 
    \begin{equation*}
        \sum_\lambda c_{\lambda^c} d_\lambda(M) = \Vol_s(P(M)),
    \end{equation*}
    where $c_{\lambda^c}$ are as in \cref{prop:VolDeg}.
\end{corollary}
\begin{proof}
    For a fixed positive integer $s \leq r(n-r)$, notice that both 
    \begin{align*}
        \Vol_s(P(-)) \colon & S(r,[n]) \to \ZZ \text{ and }\\
        \deg \big((\sum_\lambda d_\lambda(-) \sigma_\lambda )\sigma_{(1)}^s \big) \colon & S(r,[n]) \to \ZZ,
    \end{align*}
    where the sum is over all partitions $\lambda \subseteq r\times(n-r)$ with $|\lambda| = r(n-r)-s$, are valuative matroid invariants. By Theorem 6.3 of \cite{DerksenFink} all valuative matroid invariants are determined by their values on the Schubert matroids $SM_I$. Since the Schubert matroids are representable over $\CC$ we know that $\deg((\sum_\lambda d_\lambda(SM_I) \sigma_\lambda )\sigma_{(1)}^s) = \Vol_s(P(SM_I))$ by \cref{prop:VolDeg}. Hence these invariants are equal for any matroid. The corollary follows from setting $s = n-\kappa(M)$.
\end{proof}

We will use this to compute the Schubert coefficients of sparse paving matroids. In light of the previous section we may assume that all matroids are connected. First we compute The Schubert coefficients of the minimal matroids.
\begin{lemma}
    \label{lem:minimalMatroid}
    The Schubert coefficients of the minimal matroid $T_{r,n}$ are
    $d_\lambda(T_{r,n}) = 0$ for $\lambda \neq h^c$ and $d_{h^c}(T_{r,n}) = \beta(T_{r,n}) = 1$, i.e.\ $\Sc(T_{r,n}) = \sigma_{h^c}$.
\end{lemma}

\begin{proof}
    By \cref{prop:VolDeg}
    \begin{equation}
        \label{eq:minimalMatroidEq}
        \sum_\lambda c_{\lambda^c} d_\lambda(T_{r,n}) = \Vol(P(T_{r,n})).
    \end{equation}
    The matroid polytope $P(T_{r,n})$ is a pyramid of lattice height $1$ over the product of simplices $\Delta_{r-1} \times \Delta_{n-r-1}$ and hence has volume $\Vol(P(T_{r,n})) = \binom{n-2}{r-1}$; see also example 7.4 of \cite{FerSch}. Next we compute the coefficient of the hook partition $c_h$. By using Pieri's formula and considering the expansion of $\sigma_{(1)}^{n-1}$ in $A^\bullet(G(r,n))$ we see that $c_{\lambda^c} = |SYT(\lambda^c)|$ is the number of standard Young tableaux with shape $\lambda^c$. The following formula for $|SYT(\lambda^c)|$ is given in Theorem $1$ of \cite{FrameRobinsonThrall}\footnote{This paper gives the formula for the degree of the corresponding irreducible representation of the symmetric group on $|\lambda|$ elements. There is a well known basis of this representation in bijection with standard Young tableaux of shape $\lambda$.}:
    \begin{equation*}
        |SYT(\lambda^c)| = \frac{|\lambda^c|!}{\prod_{b\in \lambda^c} hl(b)}
    \end{equation*}
    where the product is over all boxes in the Young diagram of $\lambda^c$ and $hl(b)$ is the hook length of $b$. Applying this formula to the hook partition $h$ gives
    \begin{align*}
        c_h = |SYT(h)| &= \frac{(n-1)!}{(n-1)(r-1)!(n-r-1)!}\\
        &= \frac{(n-2)!}{(r-1)!(n-2 -(r-1))!}\\
        &= \binom{n-2}{r-1}
    \end{align*}
    hence $\Vol(P(T_{k,n})) = c_h$. Since $d_{h^c}(T_{r,n}) = \beta(T_{r,n}) = 1$ \cref{eq:minimalMatroidEq} becomes
    \begin{equation*}
        c_h + \sum_{\lambda\neq h^c} c_{\lambda^c}d_\lambda(T_{r,n}) = \Vol(P(T_{r,n})).
    \end{equation*}
    The integers $c_{\lambda^c}$ are positive since it is the degree of the Schubert variety $X_\lambda(V_\bullet)$ and the Schubert coefficients $d_\lambda(T_{r,n})$ are non-negative since $T_{r,n}$ is representable over $\CC$. Hence $d_\lambda(T_{r,n}) = 0$ for all partitions $\lambda \neq h^c$.
\end{proof}

\begin{remark}
\label{rmk:minimal}
    In light of \cref{lem:minimalMatroid} we might guess that if \(x \in G(r,n)\) represents the minimal matroid \(T_{r,n}\), then the torus orbit closure \(\T\) is equal to the Schubert variety \(X_{h^c}(V_\bullet)\). Indeed one can see that this holds if the $r$-dimensional part of \(V_\bullet\) is the vector space spanned by the standard vectors denoting the non-parallel columns of a matrix representing \(x\).
\end{remark}

By using \cref{lem:minimalMatroid} we are now ready to prove \cref{thm:1}.
\begin{proof}[Proof of \cref{thm:1}]
    Since $M$ is sparse paving there is a matroid subdivision of $P(U_{r,n})$ containing as facets $P(M)$ and $k$ polytopes of matroids isomorphic to $T_{r,n}$. Since $d_\lambda$ are additive and valuative matroid invariants this subdivision yields
    \begin{equation*}
        d_\lambda(U_{r,n}) = d_\lambda(M) + k\cdot d_\lambda(T_{r,n}) \iff d_\lambda(M)  = d_\lambda(U_{k,n}) - k\cdot d_\lambda(T_{r,n}).
    \end{equation*}
    By \cref{lem:minimalMatroid} if $\lambda \neq h^c$ then $d_\lambda(M) = d_\lambda(U_{r,n})$ and $d_{h^c}(M) = \beta(M) = \beta(U_{r,n}) - k$.
\end{proof}
Now \cref{cor:1} follows since $U_{r,n}$ is representable over $\CC$ so $d_\lambda(M) = d_\lambda(U_{r,n})$ and $\beta(M)$ are non-negative.

The proof of \cref{thm:1} applies to the following more general case.
\begin{corollary}
    Let $M$ and $M'$ be connected matroids of rank $r$ on $[n]$ and assume there is a matroid subdivision of $P(M')$ containing as facets $P(M)$ and $k$ polytopes of matroids isomorphic to $T_{r,n}$. Then $d_\lambda(M) = d_\lambda(M')$ for all $\lambda \neq h^c$.
\end{corollary}
\begin{example}
    Some well known examples of matroids that are not representable over $\CC$ are the Fano matroid $F$, the non-Pappus matroid $NP$ and the Vamos matroid $V$. These are sparse paving matroids and the classes of $F$, $NP$ and $V$ in $A^\bullet(G(r,n))$ are shown below.
    \begin{align*}
        \Sc(F) = & 6 \sigma_{(4,2,0)} + 3 \sigma_{(4,1,1)} + 3 \sigma_{(3,3,0)} + 8 \sigma_{(3,2,1)} + \sigma_{(2,2,2)} \in A^\bullet(G(3,7))\\
        \Sc(NP) = & 15 \sigma_{(6,4,0)} + 15 \sigma_{(6,3,1)} + 6 \sigma_{(6,2,2)} + 13 \sigma_{(5,5,0)} \\
        & + 24 \sigma_{(5,4,1)} + 15 \sigma_{(5,3,2)}. + 6 \sigma_{(4,4,2)} + 3 \sigma_{(4,3,3)} \in A^\bullet(G(3,9))\\
        \Sc(V) = & 4 \sigma_{(4,4,1,0)} + 20 \sigma_{(4,3,2,0)} + 12 \sigma_{(4,3,1,1)} + 12 \sigma_{(4,2,2,1)} \\
        & + 15 \sigma_{(3,3,3,0)} + 20 \sigma_{(3,3,2,1)} + 4 \sigma_{(3,2,2,2)} \in A^\bullet(G(4,8))
    \end{align*}
\end{example}

One might try to expand this argument to prove non-negativity of the Schubert coefficients for a class of matroids larger than the sparse paving matroids. In what follows we give an example showing this does not work for paving matroids that are not sparse paving.

As we have seen, sparse paving matroids can be described as those matroids obtained from uniform matroids by cutting of minimal matroid polytopes from the hypersimplex. In a similar way, as described in \cite{HMMMNVYPanhandle}, one can describe paving matroids as those obtained from $U_{r,n}$ by cutting of panhandle matroids. Recall that the panhandle matroid $\Pan_{r,r,n}$ is the minimal matroid $T_{r,n}$ and $\Pan_{r,n-1,n}$ is the uniform matroid $U_{r,n}$. The next example shows that \cref{thm:1} does not hold for paving matroids that are not sparse paving.
%If the argument in \cref{thm:1} is to be extended to paving matroids we would need the Schubert coefficients of panhandle matroids to be zero, except for $d_{h^c}(\Pan_{r,s,n})$. The next example shows that this is not the case. \todo{change this text now that i have changed the example?}

\begin{example}
    The smallest non-minimal, non-uniform example of a panhandle matroid is $\Pan_{2,3,5}$. There is a matroid subdivision of $U_{2,5}$ of the form
    \begin{equation}
        \label{eq:PanMinUni}
        P(U_{2,5}) = P(\Pan_{2,3,5}) \cup P(T_{2,5}).
    \end{equation}
    That is $\Pan_{2,3,5}$ is a sparse paving matroid, and $T_{2,5}$ is a paving matroid that is not sparse paving. 
    By \cref{lem:minimalMatroid} and \cref{thm:Klyachko} we have
    \begin{equation*}
        d_{(1,1)}(T_{2,5}) = 0 \neq 1 = d_{(1,1)}(U_{2,5})
    \end{equation*}
    so \cref{thm:1} does not hold for paving matroids in general. For completeness by \cref{thm:Klyachko}, \cref{lem:minimalMatroid} and \cref{eq:PanMinUni} we have
    \begin{align*}
        \Sc(U_{2,5}) &= \sigma_{(1,1)} + 3 \sigma_{(2)}, \\
        \Sc(T_{2,5}) &= \hspace{1.5cm} \sigma_{(2)} \text{ and }\\
        \Sc(\Pan_{2,3,5}) &= \sigma_{(1,1)} + 2\sigma_{(2)}.
    \end{align*}
\end{example}
This leads us to state the following conjecture.
\begin{conj}
    Let $M$ be a connected matroid of rank $r$ on $[n]$ then $M$ is sparse paving if and only if $d_\lambda(M) = d_\lambda(U_{r,n})$ for all partitions $\lambda \neq h^c$.
\end{conj}
